\begin{table*}[t]
\caption{The playful health technology (PHT) framework questions, with examples for each answer option.}
\label{tab:pht}
\small
\begin{tabular}{p{0.2\textwidth}p{0.72\textwidth}}
\toprule
Framework question & Answer options \\
\midrule
Fun vs. Utility & (1) totally fun first (i.e., Entertainment game reappropriated) \newline (2) mostly fun first (i.e., design workshop with disabled stakeholders) \newline (3) both fun and utility first (i.e., design workshop with diverse stakeholders) \newline (4) mostly utility first (i.e., primary utilitarian mechanic designed first) \newline (5) totally utility first (i.e., made by clinicians or specialists, fun incorporated later) \\ \addlinespace
Play vs. Game & (1) unstructured play (i.e., no rules, making art) \newline (2) semi-structured play (i.e., playground activities) \newline (3) flexible structure with rules (i.e., improv) \newline (4) flexible game (i.e., board game with house rules) \newline (5) game with rigid rules (i.e., video game) \\ \addlinespace
Skill vs. Chance & (1) entirely skill-based (i.e., trivia, sports) \newline (2) mostly skill-based (i.e., mario kart) \newline (3) equally skill and chance-based (i.e., catan) \newline (4) mostly chance-based (i.e., Uno) \newline (5) entirely chance-based (i.e., all dice) \\ \addlinespace
Solo vs. Social & (1) entirely solo (i.e., solitaire) \newline (2) mostly solo (i.e., playing against AI in single-player StarCraft) \newline (3) mix of solo and social affordances (i.e., animal crossing) \newline (4) mostly social, but aspects are independent (i.e., house on hill haunt transition) \newline (5) entirely social (i.e., tag) \\ \addlinespace
Sequential vs. Simultaneous & (1) entirely turn-based (i.e., tic tac toe) \newline (2) follows a set of steps (i.e., viticulture's dynamic turn ordering) \newline (3) turns are taken, but some actions can be taken at any time (i.e., pandemic) \newline (4) most actions can be taken at any time, but there are some phases (i.e., PvP death reset timer) \newline (5) entirely simultaneous (i.e., race) \\ \addlinespace
Synchronous vs. Asynchronous & (1) entirely synchronous (i.e., real time strategy game) \newline (2) mostly synchronous with a few asynchronous affordances (i.e., league of legends) \newline (3) equal mix of synchronous and asynchronous affordances (i.e., helldivers 2) \newline (4) mostly asynchronous with few synchronous affordances (i.e., cookie clicker) \newline (5) entirely asynchronous (i.e., chess by postage mail) \\ \addlinespace
Competitive vs. Collaborative & (1) entirely competitive (i.e., spit card game) \newline (2) mostly competitive, but sometimes collaboration is important (i.e., forbidden island) \newline (3) mix of competitive and collaborative (i.e., mario party 2 vs 2 mini-games) \newline (4) mostly collaborative (i.e., animal crossing) \newline (5) entirely collaborative (i.e., pandemic) \\ \addlinespace
Symmetrical vs. Asymmetrical & (1) entirely symmetrical (i.e., tic tac toe) \newline (2) mostly symmetrical (i.e., undertale's differing NPC behaviors based on previous playthroughs) \newline (3) both symmetrical and asymmetrical (i.e., pandemic with its similar turn structure but different character abilities) \newline (4) mostly asymmetrical (i.e., tag) \newline (5) entirely asymmetrical (i.e., keep talking, nobody explodes) \\ \addlinespace
\bottomrule
\end{tabular}
\end{table*}
